\section{Traffic engineering with segment routing}

In this section we discuss the variants of Problems \ref{prob:tep} and \ref{prob:tep-mul} when
instead of using arbitrary paths for routing, we use sr-paths. The SR version of the TE problem is
obtained by replacing paths by sr-paths, $\cnt(p_i, e)$ by $\r(\sr{p}_i, e)$ and adding a constraint on the segment cost of 
the sr-paths found. Recall that $\r(\sr{p}_i, e)$ was defined on Chapter \ref{chapter:sr-optimal} and represents the 
proportion of the traffic that traverses edge $e$ when routing over sr-path $\sr{p}$.

This last constraint is important to ensure that the paths can be supported by the routers in the network.
This leads to the definition of the following two problems which, respectively, correspond to Problems \ref{prob:tep} and \ref{prob:tep-mul}.

\begin{problem}{Segment routing traffic engineering}
\label{prob:srte}
\textbf{Input:} A network $G$, a set of demands $\mathcal{D} = \{d_1, \ldots, d_n\}$ on $G$ and $k \in \mathbb{N}$.

\textbf{Output:} The minimum factor $\lambda \geq 0$ and set of sr-paths on $G$, $\sr{p}_1, \ldots, \sr{p}_n$ such that $\sr{p}_i$ is a 
sr-path from $s_i$ to $t_i$ with $\cost(\sr{p}_i) \leq k $ and for each link $e \in E(G)$ it holds
$$
\sum_{i = 1}^n \r(\sr{p}_i, e) \cdot \vol(d_i) \leq \lambda \cdot \bnd(e).
$$ 
\end{problem}

\begin{problem}{Segment routing traffic engineering}
\label{prob:srte-mul}
\textbf{Input:} A network $G$, a set of demands $\mathcal{D} = \{d_1, \ldots, d_n\}$ on $G$ and $k \in  \mathbb{N}$.

\textbf{Output:} The minimum factor $\lambda \geq 0$ and set of sr-paths on $G$, $\sr{\mathcal{P}}_1, \ldots, \sr{\mathcal{P}}_n$ such that
$\sr{\mathcal{P}}_i$ is a 
non-empty set of pairs $(\sr{p}, f)$ where $\sr{p}$ sr-path from $s_i$ to $t_i$ with $\cost(\sr{p}_i) \leq k $, $f \in [0, 1]$ and for each link $e \in E(G)$ it holds
$$
\sum_{i = 1}^n \sum_{(\sr{p}, f) \in \sr{\mathcal{P}}_i} f \cdot \r(\sr{p}, e) \cdot \vol(d_i) \leq \lambda \cdot \bnd(e).
$$ 
\end{problem}

The minimum segmentation algorithm provides a way to translate solutions of Problems \ref{prob:tep} and \ref{prob:tep-mul}
into solutions of Problems \ref{prob:srte} and \ref{prob:srte-mul}, respectively. The advantage of this approach is that
makes it possible to leverage existing algorithms for solving these widely studied graph problems. The drawback of course
is that since these algorithms are oblivious to segment routing, there is no way of knowing whether or not the output paths will
require too many segments for routers to be able to support them.

We first evaluate the segment cost mentioned above. 
We already did this in the previous chapter when we considered the problem of computing minimum latency sr-paths.
This is a way to evaluate how necessary it is to develop dedicated algorithms and how often we can get away by
simply segmenting graph centric solutions using the minimum segmentation algorithm proposed in Chapter \ref{chapter:sr}.
Perhaps in the future routers will be able to support a high amount of segments
and whenever this is the case, it will become fruitless to develop dedicated algorithms.

% We used the dataset from the Repetita project to evaluate the segment cost of solutions to Problem \ref{prob:tep}. For each topology
% we executed Algorithm \ref{algo:mcf-seg} defined in the previous section and segment all the paths from the solution. Figure 
% \ref{fig:minseg_MCF} shows a CDF  of the maximum segment cost of minimum segmentation the path sets correspond to each demands,
% over all topologies.
% 
% \begin{figure}
% \begin{center}
% \includegraphics[scale=0.6]{./Network-lib/data/plot/todo.pdf}
% \end{center}
% \caption{CDF of the segment cost of the minimal segmentations of solution computed with Algorithm \ref{algo:mcf-seg}.}
% \label{fig:minseg_MCF}
% \end{figure}


\subsection{Existing MIP models and algorithms for SRTE}

We now have all the tools needed to describe existing models for the traffic engineering problem with segment routing.
The first approach that was developed by Bathia et al and considers only sr-paths of the form $\langle s, x, t \rangle$ to route a demand from
$s$ to $t$ \cite{bhatia}. That is, it restricts the set of admissible sr-paths to sr-paths with a single detours towards a given
router $x$ on the network. The main advantage of doing this is that we obtain a very efficient model since for each demand 
there is a single decision to make: which intermediate router to use.

% To simplify the notations and model formulations, it is convenient to define the set of all sr-paths with a segment cost lower
% than a given constant.
% 

% 
% With this notation,

Let $\PB = \left\{ \langle \src(d), x, \dst(d) \rangle \mid x \in V(G) \right\}$. We define binary variables $x_{dp}$ such that $x_{dp} = 1$ if and only if demand $d$ is routed
over sr-path $p$.  Their model can be formulated as follows.

\begin{center}
\begin{tabular}{crcllr}
\multicolumn{5}{l}{$\bhatia(G, \mathcal{D})$} \\[0.5cm] 
$\displaystyle \mathbf{min}$ & $\lambda$ & & & & \\[0.5cm]
$\textbf{s.t.}$ & $\displaystyle \sum_{p \in \PB} x_{d p}$ & $=$ & $1$ & $\forall d \in \mathcal{D}$ \\[0.5cm] 
& $\displaystyle \sum_{\sr{p} \in \PB} \sum_{d \in \mathcal{D}} \r(\sr{p}, e)  \cdot \vol(d) \cdot x_{dp}$ & $\leq$ & $\lambda \cdot \bnd(e)$ & $\forall e \in E(G)$ \\[0.5cm]
& $x_{dp} \in \{ 0, 1 \}$ & & & $\forall p \in \PB, \ \forall d \in \mathcal{D}$
\end{tabular}
\end{center}


The first set of constraints specifies that each demand has to be served by exactly one path. The second set of constraints ensures that no edge 
carries more than $\lambda \cdot \bnd(e)$ traffic.

The model that we use for our column generation solution is a generalization of this one where we replace the set of paths considered, $\PB$
by the set of all sr-paths with at most a given segment cost, $\mathcal{\sr{P}}_k(G)$. 
Renaud Hartert already had proposed this generalization in his thesis \cite{hartert18}. The only difference between his model and ours is that we also
consider adjacency segments in our sr-paths whereas his model only considered node segments. We discuss our model in more detail in the next section.

In his thesis, Hartert also proposed another model called \emph{segment model}. This model is closely related to the IP version of the model that we presented for the
multi-commodity flow \mcflp. In the integral MCF model, we use variables $x_{ed}$ to indicate whether demand $d$ is routed though edge $e$. Then
we use classic flow conservation constraints to ensure that if $x_{ed} = 1$ then we have a path $(e_1, \ldots, e_l)$ starting at
$\src(d)$ and ending at $\src(d)$ such that $x_{e_1 d} = 1, \ldots, x_{e_l d} = 1$. In the segment model, instead of edges we consider the shortest
path subgraphs. Concretely, we use variables $x^d_{u v}$ for $d \in \mathcal{D}, u, v \in V(G)$
defined such that $x^d_{u v} = 1$ if and only demand $d$ is routed over the shortest paths between nodes $u$ and $v$, that is, the subgraph $\sp(u, v)$. 
Then we use the same flow conservation constraints to ensure that whenever $x^d_{uv} = 1$ then there exists a sequence of nodes
$(v_1 = s, v_2, \ldots, v_k = t)$ such that $x^d_{v_1 v_2} = x^d_{v_2 v_3} = \ldots = x^d_{v_{k - 1} v_k} = 1$ meaning that we route the demand
$d$ by following the shortest paths from $v_1$ to $v_2$, then from $v_2$ to $v_3$ and so on. In other words, we use sr-path 
$\langle v_1, v_2, \ldots, v_k \rangle$ to route $d$.

\begin{center}
\begin{tabular}{rcllr}
\multicolumn{5}{l}{$\srteseg(G, \mathcal{D})$} \\[0.5cm] 
\multicolumn{3}{l}{$\mathbf{min} \quad \lambda$} & $\textbf{s.t.}$ & \\[0.5cm]
$\displaystyle \sum_{d = 1}^r \sum_{u, v \in V(G) : u \neq v} x^d_{uv} \cdot \r(u, v, e) \cdot \vol(d)$ & $\leq$ & $\lambda \cdot \bnd(e)$ & $\forall e \in E(G)$ & \\[0.5cm]
$\displaystyle \sum_{u \in V(G) \setminus \{ v \}} x^d_{uv} - \sum_{u \in V(G) \setminus \{ v \}} x^d_{vu}$ & $=$    &  $0$ & $\forall d = (s, t, \nu) \in \mathcal{D}$, & \\[-0.2cm]
& & & $\forall v \in V(G) \setminus \{ s, t \}$ & \\[0.5cm]
$\displaystyle \sum_{u \in V(G) \setminus \{ s \}} x^d_{us} - \sum_{u \in V(G) \setminus \{ s \}} x^d_{us}$ & $=$    & $-1$ & $\forall i \in \{ 1, \ldots, r \}$, \\[-0.2cm]
& & & $\forall d = (s, t, \nu) \in \mathcal{D}$ \\[0.5cm]
$\displaystyle \sum_{u \in V(G) \setminus \{ t \}} x^d_{ut} - \sum_{u \in V(G) \setminus \{ t \}} x^d_{ut}$ & $=$    &  $1$ & $\forall i \in \{ 1, \ldots, r \}$, \\[-0.2cm]
& & & $\forall d = (s, t, \nu) \in \mathcal{D}$ \\[0.5cm] 
$x^d_{uv}$  &    $\in$    &  $\{0, 1\}$  & $\forall e \in E(G), \ \forall d \in \mathcal{D}$ & \\[0.5cm]
$\lambda \geq 0$ & & & &    
\end{tabular}
\end{center}


As it is, this model does not make any guarantees on the number of segments used to route a demand. We can overcome this by adding the following additional constraint
limiting the maximum number of node segments in the sr-paths to $k$:

$$
\sum_{u, v \in V(G) : u \neq v} x^d_{u, v} \leq k \quad \forall d \in \mathcal{D}.
$$

In all these models, by relaxing the integrability constraints to allow the variables to take any real value in $[0, 1]$ we automatically get a polynomial time solvable
LP for Problem \ref{prob:srte-mul} (over the restricted set of sr-paths that each model considers). 

\subsubsection*{Model comparison}

Due to its simplicity, optimal solutions of $\bhatia$ can be computed very efficiently. It it a relaxation of Problem \ref{prob:srte} because it significantly restricts the
sets of sr-paths that can be used by considering only paths in $P_1$. The generalization of this model consisting of replacing $P_1$ by $\mathcal{\sr{P}}_k$ removes this 
restriction by considering any possible sr-path whose segment cost is at most $k$ to route any given demand. The problem of this model is that it contains a number of variables
which is exponential with respect to $k$ since $|\Pk| = O(|G|^k)$. This means that even for small values of $k$ it will not scale. Hartert showed in his thesis that even for $k \geq 3$ the number of 
variables is too big to allow good solutions to be found in a reasonable amount of computation time. We will show how we use column generation to overcome this problem the next section.

The segment model presented above overcomes this problem by implicitly representing sr-paths with flow conservation constraints. With this, it manages
to model sr-paths of segment cost up to $k$ using only $|V(G)|^2 \cdot |\mathcal{D}|$ variables. One drawback of this solution is that it
cannot represent sr-paths containing adjacency segments. Moreover, it turns out that on the larger topologies this model quickly becomes too large as well. This is not
surprising since with one demand per pair of nodes the total number of variables is about $|V(G)|^4$.

Other researchers have proposed heuristic techniques for solving Problem \ref{prob:srte}. Hartert et al. proposed in DEFO \cite{defo,hartert2015solving} a Large
Neighborhood Search technique combined with Constraint Programming. Later, Gay et al. proposed in \cite{steven} to use standard local search to iteratively improve the current solution. 
These heuristic approaches are rather efficient but provide no way of knowing how far from the optimal value they end up.
Moreover, all the existing approaches only consider node segments in their formulation. They are thus not able to fully exploit the flexibility of segment routing with adjacency segments.
Not only that but we have shown in our experiments that, on instances where only a few routing configuration lead to good solution, these heuristic approaches have difficulty of finding them. 
This is a common drawback of LS since in such cases it becomes very unlikely for the search to be able to reach these very precise configurations.
